\documentclass[11pt]{article}

\usepackage[a4paper,margin=1in]{geometry}
\usepackage{amsmath}
\usepackage{booktabs}
\usepackage{hyperref}

\hypersetup{colorlinks=true,linkcolor=blue,citecolor=blue,urlcolor=blue}

\title{Evidence-Aligned Repetition-Code Benchmarking:\\
Corrected Decoder Baselines, Structured Noise Tests, and Hardware Execution Checks}
\author{Project Manuscript Draft}
\date{\today}

\begin{document}
\maketitle

\begin{abstract}
This draft reports a corrected repetition-code benchmarking workflow. The strongest supported simulation results are (i) distance scaling under Stim detector-model decoding and (ii) weighted versus unweighted MWPM comparison on identical shot samples. Earlier over-broad interpretations are removed: detector-event voting is treated as a weak control, not a proper codeword-majority baseline, and heterogeneous/drift claims are tied only to corrected per-qubit/per-round implementations. Hardware results are presented as constrained execution checks unless repeated-syndrome decoding evidence is explicitly available from saved records.
\end{abstract}

\section{Scope and Claim Policy}
This manuscript follows a strict evidence policy:
\begin{enumerate}
    \item Every quantitative claim must map to saved artifacts generated by the current codebase.
    \item Simulation and hardware claims are separated; hardware is not used to imply threshold-level correction.
    \item Confidence statements are based on integer failures and shot counts, not rounded rates alone.
\end{enumerate}

\section{Corrected Methods}
\subsection{Decoder-comparison correction}
The decoder-comparison pipeline now evaluates all methods on \emph{identical sampled shots}. A single measurement sample is converted into detector events and observable flips, then reused across decoders.

The method labeled \texttt{majority\_vote} in earlier drafts was corrected. In the current code:
\begin{itemize}
    \item \texttt{final\_data\_majority} is the proper final-data baseline.
    \item \texttt{detector\_count\_heuristic} is a weak detector-event control.
    \item \texttt{unweighted\_mwpm} and \texttt{mwpm} are compared with pairwise win/loss counts.
\end{itemize}

\subsection{Bias, heterogeneity, and drift corrections}
Bias experiments use explicit round-aware circuit construction with data channels applied before syndrome/readout. Heterogeneous experiments now apply per-qubit rates directly (same-mean control versus local-defect scenarios). Temporal-drift experiments now apply explicit per-round schedules (constant, split low/high schedule, and measurement-only perturbation) instead of collapsing schedules to a single mean.

\subsection{Hardware correction path}
The corrected hardware stream uses repetition circuits with repeated syndrome extraction (data + ancilla qubits), explicit delays, syndrome-aware logical decoding from returned measurement bitstring records/counts, and Wilson confidence intervals. These hardware outputs are interpreted as constrained workflow evidence unless broader validation conditions are met.

\section{Results Supported by Current Artifacts}
\subsection{Primary supported result: simulation decoder behavior}
Distance scaling and weighted-vs-unweighted MWPM remain the primary supported findings. A corrected high-statistics slice at $d=11$, $p=0.02$ (100{,}000 shots per seed) produced lower failure counts for DEM-weighted MWPM than unweighted MWPM across all recorded seeds, with integer-count manifests saved in project run manifests.

\subsection{Secondary supported result: corrected structured-noise pipelines}
Corrected heterogeneity and temporal-schedule pipelines execute end-to-end with explicit scenario metadata, same-sample decoder comparisons, and reproducible seeds. These outputs are suitable for technical review of method behavior, not yet broad generalization.

\subsection{Hardware interpretation}
A corrected cloud run on \texttt{ibm\_fez} (job \texttt{db0ph0ivog1s73fh5km0}, 1200 shots/spec) completed and produced syndrome-validation artifacts. The encoded variants achieved logical-success estimates of 1.000 (r=1, logical 0), 0.990 (r=1, logical 1), 1.000 (r=3, logical 0), and 0.970 (r=3, logical 1), while the duration-matched unencoded controls were 0.9925 (logical 0) and 0.9000 (logical 1), with Wilson 95\% intervals recorded per row in saved CSV artifacts. These results provide constrained syndrome-decoding checks under the implemented protocol; they do not, by themselves, establish fault-tolerant hardware advantage.

\section{Limitations}
\begin{itemize}
    \item Repetition code benchmarks a restricted error-correction setting.
    \item Some structured-noise and hardware conclusions remain conditional on run budgets, queue conditions, and backend state.
    \item No threshold, universality, or device-independent advantage claim is made.
\end{itemize}

\section{Conclusion}
After correction, the project contributes a clearer and more defensible benchmark workflow: identical-shot decoder comparisons, explicit structured-noise implementations, and constrained hardware execution checks with confidence accounting. The strongest claims are simulation-based and decoder-centric; hardware claims are intentionally scoped.

\end{document}
